\chapter{Operationeel profiel}


Om het onderzoek wat in dit verslag is beschreven te onderbouwen, is het van belang om het operationeel profiel van transportschip Jansje te bepalen op een manier waar de eigenaar, Wim van Baalen, op kan vertrouwen. Dit sluit aan bij de eerdergenoemde deelvraag 1:

\begin{quote}
    \textit{\textbf{Hoe wordt het operationeel profiel van Jansje op een representatieve en 
    reproduceerbare wijze vastgesteld?}}
\end{quote}

\noindent Om een geschikte, duurzame aandrijflijn voor Jansje te kunnen ontwerpen, moet eerst worden vastgesteld hoe het schip daadwerkelijk wordt gebruikt. Het operationeel profiel beschrijft daarom de frequentie, duur, snelheid, vermogensvraag, hotel-load, financiële randvoorwaarden en meer van het schip. Deze gegevens vormen de basis van het onderzoek naar energieopslagsystemen. Hiervoor wordt eerst de onderzoeksmethode toegelicht. Vervolgens worden het vaarprofiel, het vermogen en energieverbruik, het stilliggen en de laadmogelijkheden, en de belangrijkste randvoorwaarden behandeld.


\section{Onderzoeksmethode}

Om het operationeel profiel te bepalen is een interview gehouden met de eigenaar van Jansje. In combinatie met eerder mailcontact en dit interview is alle benodigde informatie verkregen. Omdat een deel van de gevraagde informatie echter niet direct beschikbaar was, is het operationeel profiel aanvullend onderbouwd aan de hand van vergelijkbare vaartuigen die worden beheerd in de haven van Maassluis door Loods M. De vaartijden van deze schepen zijn openbaar beschikbaar en bieden een representatief referentiekader voor het typische gebruikspatroon van historische vaartuigen in dit vaargebied.

\section{Vaarprofiel}
Het vaarprofiel is bepaald door het interview met de eigenaar. \\
\noindent Uit het inverview \cite{vanbaalen2026interview} zijn de volgende conclusies getrokken:
\begin{itemize}
    \item Het schip vaart tussen mei en oktober.
    \item Het schip vaart ongeveer 30 dagen per jaar.
    \item Een vaartrip duurt gemiddeld 6 uur en maximaal 8 uur.
    \item De gemiddelde vaarsnelheid is 10-12 km/h.
    \item De maximale vaarsnelheid is 14 km/h.
    \item Er is vooral piekvermogen nodig tijdens het wegvaren en aanmeren, en manoeuvreren. Hoevaak dit gedaan moet worden per vaart is voornamelijk afhankelijk van het aantal bruggen dat gepasseerd moet worden.
    \item De routes zijn willekeurig. 
    \item De boot gaat naar events waar ze uitgenodigd worden, van Amsterdam tot Dordrecht, overal in Nederland.
\end{itemize}


\section{Vermogen en energieverbruik}
\subsection{Piekvermogen en manoeuvreren}

De boot wordt op dit moment aangedreven door een motor van 100pK van General Motors \cite{Jansje}. Omgerekend is dat bijna 80 kW. Verder is de hotel-load bestaat uit alle elektriciteitconsumerende apparaten aan boord die geen onderdeel zijn van de aandrijflijn.
Uit het interview en een rondleiding is de conclusie getrokken dat hier rekening mee gehouden moet worden.
\noindent Uit het interview bleek ook dat er bepaalde apparaten aanwezig zijn op het schip die vallen onder de Hotel load:
\begin{itemize}
    \item Verlichting
    \item Koelkast
    \item Navigatie
    \item Koken
    \item Boiler
\end{itemize}

\section{Stilliggen en laadmogelijkheden}
Tijdens het interview is naar voren gekomen dat met transportschip Jansje ook naar events wordt gevaren verder in het land. Voor bijvoorbeeld SAIL in Amsterdam, wordt dan twee dagen achtereenvolgend gevaren met Jansje. De eigenaar van Jansje heeft aangegeven dat de vaartijd van acht uur per dag een goede indicatie is, en dat ervan uitgegaan kan worden dat er een elektriciteits aansluiting van 230V beschikbaar is als het schip ligt aangemeerd.\\



\section{Resultaten}

In deze paragraaf zijn de parameters van het operationeel profiel uitgezet:

\begin{table}[H]
    \centering
    \caption{Geregistreerde vaartijden van referentieschepen beheerd door Loods M te Maassluis}
    \label{tab:vaartijden}
    \begin{tabular}{llll}
        \toprule
        \midrule
        \textbf{Vaarprofiel \& routes} & \\
        \midrule
        \bottomrule
        Hoeveel vaart jansje per jaar? & 30 dagen \\
        Hoeveel uur vaart Jansje gemiddeld? & 6-8 uur \\
        wat is de gemiddelde vaarsnelheid? & 10-12 km/h \\
        Welke afstand vaart jansje gemiddeld? & 36-80 km\\
        
        \toprule
        \midrule
        \textbf{Vermogen \& energieverbruik} & \\
        \midrule
        \bottomrule
        Wat is het piekvermogen van de oude motor? & 100pk (73.55 kW)\\
        Hoeveel piekvermogen is er nodig tijdens het varen? & ??\\
        Wat is de geschatte hotel-load? & ??\\
        
        \toprule
        \midrule
        \textbf{Stilligen \& Laadmogelijkheden} & \\
        \midrule
        \bottomrule
        Hoelang ligt het schip aan wal? & ??\\
        op welke krachtstroom kan geladen worden aan wal? & ??\\
        \toprule
        \midrule
        

        
    \end{tabular}
\end{table}




\section{Conclusie}

Met deze resultaten kan de conclusie worden getrokken dat het operationeel profiel succesvol is bepaald. Met deze resultaten kunnen de eisen, zoals later specifiek beschreven bij de verschillende deelvragen, worden opgesteld.
 
% ============================================================

